\subsection*{Abstract}
This research proposes a four-track programmatic study predicting and evaluating organized neighborhood opposition to housing development in Austin, Texas. Shifting away from omnibus associative models, the research strategy is restructured into four linked studies: (1) ex-ante multi-horizon predictive forecasting of organized opposition; (2) causal identification of the formal petition threshold via fuzzy Regression Discontinuity; (3) event-study evaluation of the HOME Initiative and H.B. 24 legislative policy shocks; and (4) a red-team governance exercise. By enforcing explicit time-stamped ``as-of'' data snapshots and establishing strict pre-deployment release gates, the study prevents predictive leakage while ensuring that municipal algorithmic models guide proactive outreach without becoming targeted exploitation tools against vulnerable neighborhoods.

\section*{1. Introduction and Background}
As the United States experiences a localized housing affordability crisis, understanding the structural barriers to urban development remains a persistent empirical challenge. Traditional models emphasize rigid municipal zoning laws, yet the procedural bottleneck is frequently driven by localized, ad-hoc political resistance, commonly framed as Not In My Backyard (NIMBY) opposition. This thesis leverages Austin, Texas, as a unique econometric laboratory. Between 2010 and 2020, rapid tech-sector growth drove severe supply constraints, inflating property values and deepening tension between neighborhood preservation and regional development.

To rigorously evaluate this friction without confounding mechanisms, this research functions as a consolidated four-track program. The empirical architecture explicitly freezes the primary analytical paths before modeling, preventing specification searching. Instead of predicting generic NIMBYism, we focus on:

\begin{itemize}
    \item \textbf{Primary Predictive Question:} At the time a zoning case is filed, what is the probability that it will face \textit{organized opposition} before final City Council action?
    \item \textbf{Primary Causal Question 1:} What is the causal friction effect of crossing the 20\% valid protest threshold on council dissent and approval behaviors under the pre-H.B. 24 regime?
    \item \textbf{Primary Causal Question 2:} How did the distinct implementation timelines of HOME Phase 1, Phase 2, and H.B. 24 fundamentally change opposition and voting dynamics for legally eligible cases?
    \item \textbf{Governance Question:} Can these early-warning predictive models be utilized for proactive municipal outreach without inadvertently creating a displacement-targeting tool for developers?
\end{itemize}

Austin operates under a council-manager system. A pro-housing majority elected in 2022 ushered in sweeping 2024 legislative reforms: expanding the HOME Act and eliminating parking requirements citywide \cite{axios2023parking}. Simultaneously, the Texas Legislature's passage of House Bill (HB) 24, taking effect September 1, 2025, restricts residential protest petitions primarily to non-comprehensive zoning changes, while maintaining the 20\% threshold in many historical contexts \cite{hamilton2024texas}. This dual evolution motivates our event-study models.

\section*{2. Literature Review}

\subsection*{2.1 Homeownership, Local Interests, and Opposition to Housing}
William Fischel's \textit{Homevoter Hypothesis} \cite{fischel2001homevoter} provides the foundational framework for understanding opposition. McCabe \cite{mccabe2016} extends this by demonstrating that homeownership itself structures preferences. Hankinson \cite{hankinson2018} shows that renters in high-cost markets also exhibit NIMBY attitudes when experiencing price anxiety. Crucially, Einstein, Palmer, and Glick \cite{einstein2019} demonstrate that public meeting participants are systematically unrepresentative, representing older, whiter homeowners. Because meeting participation is not representative, observed opposition in this study is treated distinctly from latent neighborhood preference.

\subsection*{2.2 Machine Learning, Political Behavior, and Democratic Legitimacy}
The application of machine learning to predict political behavior raises questions about democratic governance. Eubanks \cite{eubanks2018} documents how algorithmic systems often intensify the surveillance of marginalized populations. Barocas and Selbst \cite{barocas2016} show how proxies like property values can produce discriminatory outcomes. This study ensures that explicit deployment architecture addresses these procedural limits, building on Zarsky's \cite{zarsky2016} roadmap for transparent algorithmic systems.

\subsection*{2.4 Causal Inference under Distribution Shift}
To separate correlation from causality, this thesis draws from the event-study literature to rigorously test legislative shocks \cite{callaway2021difference}. Rather than relying on simple difference-in-differences, explicit date-driven staggered adoption parameters are applied. Additionally, to test extrapolation stability against concept drift over time, Out-of-Distribution methods such as Variance Risk Extrapolation (V-REx) \cite{krueger2021} and Bayesian Invariant Prediction \cite{wu2025bayesian} are implemented. However, they are treated strictly as predictive robustness checks rather than mechanisms for deriving absolute causal mechanisms from purely associative data.

\section*{3. Mixed-Methods Data and Methodology}
Instead of relying on an unstructured cross-sectional panel that permits implicit chronological overlap, the empirical foundation necessitates a clean, time-stamped spatial panel dataset. Every record across all integrated sources automatically anchors to an explicit "as-of" timestamp, allowing the system to historically reconstruct exactly what data was mathematically knowable at each discrete forecast point. This prevents temporal leakage, ensuring that backtests utilize exclusively the Census or county appraisal limits released prior to the historical prediction application date.

\subsection*{3.1 Analytical Panel Structure}
The core tables comprising the central relational warehouse include:
\begin{enumerate}
    \item \textbf{case\_master:} Case ID, nested dates (filing, notice, petition deadline, planning commission, council, withdrawal), requested zoning changes, unit limits, target FAR, status disposition.
    \item \textbf{site\_geometry:} Extracted parcel IDs, geometries, exact acreage, frontage variables, multi-buffer layers (200, 500, 1000 ft), and overlapping administrative boundaries (Council Districts).
    \item \textbf{parcel\_buffer\_snapshot:} Annually-indexed TCAD variables specifically derived for parcels within buffer thresholds (land value, structure age, homestead exemption distributions, owner-occupancy shares).
    \item \textbf{neighborhood\_snapshot:} Tract/Block Group ACS figures explicitly constrained to their historical 5-year public release dates. Includes rent burdens, tenure splits, demographic compositions, and vulnerability overlays.
    \item \textbf{policy\_calendar:} Tracks the precise 2022 council regime change; the HOME Phase 1 (Feb 5, 2024), and HOME Phase 2 (August 16, 2024) application-acceptance start dates; and the H.B. 24 Sept 1, 2025 effective date.
\end{enumerate}

Additional supportive tables log \textbf{meeting\_events} (staff/commission recommendations), \textbf{speech\_comments} (speaker stance and segmented text), \textbf{petition\_records} (validity, signature counts), and \textbf{vote\_records} (case × councilmember dynamics). A persistent \textbf{audit\_log} continuously verifies matching QA and geographical fidelity.

\subsection*{3.2 The Primary Outcome and Censoring Directives}
The \textbf{primary predictive outcome} formally identifies "Organized Opposition before final council action." This composite outcome resolves as $y=1$ if any of the following events occur:
\begin{itemize}
    \item A valid protest petition is formally filed, or
    \item At least three distinct opposing speakers/comments are historically recorded prior to the final council disposition, or
    \item Official opposition comments exceed supporting comments by a prespecified ratio.
\end{itemize}
This captures authentic mobilized friction efficiently without falling subject to the narrow constraints of statutory petitions alone. Auxiliary tasks (y1-y6) independently track specific sub-components like petition formalization, continuous opposition intensity scores, and explicit council dissenting-votes. 

\textbf{Censoring rules} strictly enforce representation bias filtering. Cases withdrawn prior to public notice are completely discarded rather than zero-classified (as no opportunity for opposition existed). Cases fundamentally missing hearing transcripts or lacking petition-database coverage are designated strictly as unknown. In addition, administratively closed interventions are omitted entirely from the primary training sample.

\subsection*{3.3 Dual-Track Feature Libraries: Deployment vs. Research}
To prevent predictive targeting algorithms from inadvertently leveraging highly sensitive vectors, our pipeline strictly enforces bifurcated feature access.

\textbf{Deployment-Eligible Feature Set:} This explicitly restricted domain feeds operations architecture. It incorporates physical site characteristics (unit shifts, PUD/TOD flags, corner-lot flags), statutory flags based unequivocally on eligibility laws (H.B. 24 status, owner vs. city initiated), aggregated buffer economics (homestead shares, land-to-total value distributions), broad local demographic footprints (median income, household-type shares, regional renter composition), proxy political volatility (historical petition rates within the council district), and macro drivers (interest rates and vacancy rates).

\textbf{Research-Only \& Audit Feature Set:} Features deployed solely restricted to normative auditing and exploratory equity impact analytics are withheld from active live-models to avoid encoding disparate impacts \cite{barocas2016}. This sequestered data includes tract-level racial and ethnic boundaries, explicit segregation overlays, historic cross-corridor proxies, personal names, and programmatic racial-surname inferences. 

\subsection*{3.4 Formalizing the Text Pipeline}
Deep NLP embeddings are structurally inappropriate to utilize if their source temporal origins (such as council transcripts) post-date the targeted prediction date. Correspondingly, text signals are only permitted late within the analytical lifecycle. Utilizing stratified random seed sampling and explicit active learning loops across 2000 speech vectors and 500 written comments, the semantic architecture bypasses generic opaque sentence-embedding models in favor of two explicit decoding layers:
\begin{itemize}
    \item \textbf{Layer 1: Stance Classification.} Assigns Support, Oppose, or Neutral polarity markers to extracted statements.
    \item \textbf{Layer 2: Frame Classification.} Executes multi-label probabilistic clustering over dominant thematic arguments (Traffic extraction, Infrastructure Capacity, Displacement Vulnerability, Neighborhood Character, and Procedural Fairness).
\end{itemize}

\section*{4. Methodology: A 4-Track Research Program}

\subsection*{4.1 Track 1: Multi-Horizon Opposition Forecasting (Ex-Ante)}
Because timeline reconstruction determines authentic algorithmic predictive power, Track 1 executes models across four distinctly constrained forecast horizons:
\begin{enumerate}
    \item \textbf{H0 (Filing Model):} Comprises exclusively what was fully known on the literal case application date (true ex ante).
    \item \textbf{H1 (Notice Model):} Integrates H0 features alongside municipal staff notices, initial neighborhood written objections, and early drafted petition text, if submitted.
    \item \textbf{H2 (Pre-Commission Model):} Combines H1 inputs with the completed planning commission packet recommendations.
    \item \textbf{H3 (Pre-Council Model):} Adds planning commission outcomes and subsequent hearing testimony. At this maximal timeframe, we restrict target analysis strictly to evaluating pre-final-vote dynamics, heavily utilizing Layer 2 NLP features.
\end{enumerate}

Given substantial negative class-imbalance realities, standard ROC-AUC metrics generate artificially inflated confidence metrics. \textbf{Precision-Recall AUC (PR-AUC)} is the singular leading discrimination metric. Sequential algorithmic modeling advances monotonically: Baseline hierarchical logistics transition into optimized Boosted Trees (LightGBM/CatBoost), capped by robust regularization methods including late-fusion MLP configurations and Variance Risk Extrapolation (V-REx) architectures enforcing invariant bounds across spatial boundaries. Candidate model selection systematically mandates validation PR-AUC maximization precisely subjected to strict calibration slope restraints bounded globally between 0.9 and 1.1.

To structurally identify failure modes mirroring eventual live-operations, outer estimation splits are intentionally non-random. Splits enforce Rolling-Origin temporal barriers, substantive Policy-Regime holdouts (such as testing pre-2022 council blocks against initial HOME phases), Spatial boundaries, and Developer (actor-oriented) historical holdout configurations. 

\subsection*{4.2 Track 2: Causal Study 1 -- The Protest Mechanism}
Moving beyond OLS correlation to address fundamental causality, we utilize the rigid Texas 20\% statutory threshold logic regulating protested neighborhood plans to deploy a definitive continuous causal model. Recovering the explicitly continuous percentage of valid signed-area share logged across internal case documentation allows for a granular Regression Discontinuity (RD) design explicitly clustered around the pre-H.B. 24 trigger threshold. Employing localized bandwidth specifications around the 20\% line strictly insulates the structural causal impact valid protest statuses inflict exclusively upon specific timeline delays and discrete council voting structures without confounding surrounding background outrage levels. (Alternatively, Double Machine Learning/DML estimations with cross-fitting are triggered if specific dimensional extraction proves incomplete).

\subsection*{4.3 Track 3: Causal Study 2 -- Policy Shocks (Event Studies)}
To evaluate municipal up-zoning dynamics, we substitute coarse 'post-policy' dummy variables for modern continuous dynamic estimation grids targeting specific event windows explicitly via standard Callaway-Sant'Anna Average Treatment Effect estimators \cite{callaway2021difference}.  HOME Phase 1 and Phase 2 constitute independent treatments mathematically keyed precisely against valid application acceptance starts. The statewide preemption framework of H.B.24 incorporates corresponding explicit structural coding evaluating exclusively parcels fulfilling legally-derived residential density applicability restrictions prior to staggered rollouts.

\subsection*{4.4 Track 4: Red-Team Governance Exercise}
The qualitative analysis explicitly transitions from rhetorical documentation toward strict governance testing designed to proactively evaluate tool deployment failure mechanisms against civic legitimacy structures. Executing focused semi-structured scenario interfaces involving public neighborhood planning staff, affordability developers, market analysts, and housing-equity officials validates how algorithm operational access structures should be regulated. Key focal objectives analyze implicit institutional exploitation parameters mapping whether revealing algorithm output logic precipitates internal institutional trust or incites active speculative capital acceleration explicitly weaponizing these vulnerability maps against established disadvantaged neighborhoods.

\section*{5. Quantitative Results and Attribution Mechanisms}

\subsection*{5.1 Ex-Ante Diagnostics and Structural Constraints}
Overall target discrimination tracking reveals pronounced dependency tied exclusively to time-of-flight information revelation thresholds. Incorporating structural H0 benchmarks systematically demonstrates standard administrative structural data generates robust baseline probability detection parameters predicting significant conflict escalation paths. However, the exact timing mechanics highlight how adding the late-fusion text classification layers explicitly bridges final PR-AUC stability targets within the localized domain architectures encompassing specific PUD variances. Overall calibration slope boundaries strongly aligned across generalized metrics. Furthermore, analyzing fold-by-fold stability matrices highlighted massive inherent distribution variance explicitly highlighting precisely why out-of-distribution tracking metrics function as core operational checks.

Feature attribution interpretations natively map the descriptive impact mechanisms propelling distinct case forecasts via explicit standard SHAP implementations. To maximize operational reader interpretability mapping causal drivers, global descriptive ablation studies document attribution exactly in structural progressive sequence arrays: starting with constrained administrative properties (unit capacity, geometry layouts), expanding across policy limits, spatial buffers (demographics and appraisal ratios), historical friction records, explicitly integrating broader macroeconomic interest metrics, integrating late-fusion text vectors alongside advanced regularization architectures sequentially. These specific attribution trails mathematically demonstrate the explicit underlying operational features commanding model classification signals while preventing overextended interpretations assuming raw attribution intrinsically verifies localized absolute causality models. 

\subsection*{5.2 Regression Discontinuity and Legal Causal Dynamics}
Preliminary evaluations focusing exactly upon the focal causal testing boundaries arrayed around the literal 20\% Texas petition volume cutoff explicitly generated highly isolated statistical divergence metrics proving threshold crossings natively instigate structural delay vectors altering explicit terminal density outcomes. Precise bandwidth boundary derivations actively evaluate structural continuity restrictions verifying exact councilmember variance arrays corresponding with active structural legal shocks. Corresponding sequential execution across the Callaway-Sant'Anna HOME impact timelines precisely evaluates immediate local property petition volatility metrics aligned precisely across specific legislative windows revealing exact local variance responses separating legally distinct local density shocks from universal underlying regional dynamics.

\section*{6. Strict Deployment Rules and Civic Governance Protocols}

Given the explicit risk explicitly demonstrated during targeted governance scenarios wherein mapping predictive weakness functionally serves as equivalent targeting architecture identifying financially vulnerable non-vocal geographic blocks, operational deployment protocols are uniformly strict. 

\textbf{No localized public-facing analytical mapping interfaces are cleared for operational deployment.} Spatial heat-map derivatives illustrating aggregated probability trajectories inherently constitute functional private acquisitions tooling vectors accelerating capital targeting against high-vulnerability neighborhoods lacking proactive political-organization frameworks \cite{barocas2016}. To fundamentally maintain the integrity of public municipal utilization architecture, the models must operate fully sequestered within verified municipal planning boundaries. Furthermore, analytical integration natively dictates utilizing explicit vulnerability identifiers exclusively for affirmative civic community integration objectives strictly prohibiting implementation models dictating pre-emptive development restriction logic. 

Consequently, formal algorithm release gates enforce explicit operational mathematical hurdles. The H0 baseline modeling architecture must quantifiably outperform optimized pure structural logistic prevalence modeling metrics by at least 25\% integrated PR-AUC minimum performance ceilings. Additionally, explicit group-validation architecture dictates strict calibration performance ranges requiring explicitly aligned 0.8--1.2 slope requirements arrayed across respective equity-based localized boundaries while constraining universal False-Negative Rate error variance bands effectively within equivalent 10\% margins limiting structural analytical bias outcomes. If predicted friction probability metrics generate massive negative correlation clustering inherently against historical spatial-vulnerability demographic boundaries, algorithm deployment is structurally aborted preventing disparate impact utilization errors entirely.

\section*{7. Conclusion}
This four-track programmatic study radically reframes Austin rezoning friction analytics by establishing exact forecasting horizons, exact causal structures targeting protest boundaries, and precise legislation-shock boundaries systematically protected by strict analytical implementation and auditing architectures. Transitioning comprehensively from raw associative predictions toward distinct causal evaluations explicitly prevents generating misguided assumptions asserting spatial probability matrices universally replicate localized neighborhood opposition drivers. By proactively mandating stringent deployment mechanisms preventing structural real estate speculative abuse, the analytical systems robustly document institutional development barriers ensuring next-generation housing density algorithms serve as protective, proactive early-warning evaluation mechanics fundamentally preserving localized civic equity objectives.
